\documentclass[10pt,a4paper]{article}
\usepackage[margin=20mm,headheight=15pt,headsep=7mm,footskip=10mm]{geometry}
\usepackage{fontspec,amsmath,amssymb,mathtools,microtype}
\setmainfont{lmroman10-regular.otf}[BoldFont=lmroman10-bold.otf,ItalicFont=lmroman10-italic.otf,BoldItalicFont=lmroman10-bolditalic.otf]
\setsansfont{lmsans10-regular.otf}[BoldFont=lmsans10-bold.otf,ItalicFont=lmsans10-oblique.otf,BoldItalicFont=lmsans10-boldoblique.otf]
\setmonofont[Scale=.86]{lmmono10-regular.otf}[ItalicFont=lmmono10-italic.otf]
\usepackage{ragged2e}\usepackage{xcolor,booktabs,longtable,array,graphicx,enumitem,fancyhdr,needspace,titlesec,hyperref,xurl}
\definecolor{AtlasNavy}{HTML}{18334C}\definecolor{AtlasTeal}{HTML}{116D73}\definecolor{AtlasGray}{HTML}{52616D}
\hypersetup{colorlinks=true,linkcolor=AtlasTeal,urlcolor=AtlasTeal,pdfcreator={Complexity Atlas},unicode=true}
\urlstyle{same}\setlength{\parindent}{0pt}\setlength{\parskip}{4.8pt}\setlength{\emergencystretch}{3em}
\renewcommand{\baselinestretch}{1.08}\setcounter{secnumdepth}{0}\setlist{nosep,leftmargin=1.5em}
\titleformat{\section}{\Large\bfseries\sffamily\color{AtlasNavy}\raggedright}{}{0em}{}
\titleformat{\subsection}{\large\bfseries\sffamily\color{AtlasTeal}\raggedright}{}{0em}{}
\titleformat{\subsubsection}{\normalsize\bfseries\color{AtlasNavy}\raggedright}{}{0em}{}
\titlespacing*{\section}{0pt}{14pt}{6pt}\titlespacing*{\subsection}{0pt}{10pt}{5pt}
\pagestyle{fancy}\fancyhf{}\fancyhead[L]{\small\sffamily\color{AtlasGray}COMPLEXITY ATLAS}
\fancyhead[R]{\small\sffamily\color{AtlasGray}ENGLISH EDITION}\fancyfoot[C]{\small\thepage}\renewcommand{\headrulewidth}{.3pt}
\newcommand{\AtlasMath}[1]{\par\begingroup\centering\setbox0=\hbox{$\displaystyle #1$}\ifdim\wd0>1.4\linewidth\typeout{ATLAS WIDE FORMULA: \the\wd0}\fi\ifdim\wd0>\linewidth\resizebox{\linewidth}{!}{\box0}\else\box0\fi\par\endgroup}
\newcommand{\AtlasInlineMath}[1]{\begingroup\setbox0=\hbox{$#1$}\ifdim\wd0>.4\linewidth\endgroup\AtlasMath{#1}\else\endgroup$#1$\fi}
\pdfstringdefDisableCommands{\def\AtlasInlineMath#1{#1}}
\DeclareMathOperator{\polylog}{polylog}\DeclareMathOperator{\dist}{dist}\DeclareMathOperator{\diam}{diam}
\DeclareMathOperator{\rank}{rank}\DeclareMathOperator{\tr}{tr}\DeclareMathOperator{\diag}{diag}
\DeclareMathOperator{\prox}{prox}\DeclareMathOperator{\KL}{KL}\DeclareMathOperator{\Var}{Var}
\providecommand{\eps}{\varepsilon}\providecommand{\R}{\mathbb R}\providecommand{\E}{\mathbb E}
\providecommand{\mathbbm}[1]{\mathbf{#1}}\providecommand{\mathds}[1]{\mathbf{#1}}
\providecommand{\llbracket}{[\![}\providecommand{\rrbracket}{]\!]}
\begin{document}
\raggedright

{\LARGE\bfseries\sffamily\color{AtlasNavy} Bilevel optimization with multiple lower solutions: uniform perturbed PL and stochastic gradients.\par}\vspace{5pt}
\textbf{Setting ID.} \texttt{BL\_\allowbreak{}SPEC\_\allowbreak{}uniform\_\allowbreak{}penalty\_\allowbreak{}pl\_\allowbreak{}stoch}

Bilevel optimization

Consider the optimistic bilevel problem \AtlasInlineMath{Y^*(x)=\operatorname*{argmin}_{y\in\mathbb{R}^{d_y}}g(x,y)}, \AtlasInlineMath{\Phi(x)=\min_{y\in Y^*(x)}f(x,y)}, and seek stationarity of \AtlasInlineMath{\min_x\Phi(x)}. Multiple global lower-level solutions are allowed; the defining minima must be attained. Among lower-level minimizers, select one minimizing the upper-level objective.

\textbf{Status.} Bounds missing

\textbf{TASK tags.} task family: Bilevel

\textbf{MODEL tags.} growth: Uniform perturbed pl; initialization: Initial gap

\textbf{ORACLE tags.} interface family: Fresh gradient

\textbf{ALGORITHMS tags.} restriction: No additional span restriction explicitly imposed

\textbf{OUTPUT tags.} certificate: No additional output certificate required

\textbf{METRIC tags.} criterion: Reduced-objective gradient norm

\textbf{GUARANTEE tags.} aggregation: Per instance; mode: Expectation

\textbf{RESOURCE tags.} accounting: Joint coordinates retained; budget: Fixed worst-case budget; measure: N grad f; measure: N grad g

\needspace{7\baselineskip}\section{1. Setting specification}
\needspace{4\baselineskip}\subsection{Task}
Consider the optimistic bilevel problem \AtlasInlineMath{Y^*(x)=\operatorname*{argmin}_{y\in\mathbb{R}^{d_y}}g(x,y)}, \AtlasInlineMath{\Phi(x)=\min_{y\in Y^*(x)}f(x,y)}, and seek stationarity of \AtlasInlineMath{\min_x\Phi(x)}. Multiple global lower-level solutions are allowed; the defining minima must be attained. Among lower-level minimizers, select one minimizing the upper-level objective.

\needspace{4\baselineskip}\subsection{Admissible model}
For every \AtlasInlineMath{x} and \AtlasInlineMath{t\in[0,\bar t]}, the perturbed lower-level objective \AtlasInlineMath{h_t(x,y)=g(x,y)+t f(x,y)} attains its minimum and satisfies \AtlasInlineMath{\|\nabla_yh_t\|^2\ge2\mu_{\mathrm{PL}}(h_t-\min_yh_t)} with the same constant \AtlasInlineMath{\mu_{\mathrm{PL}}>0}.

The joint gradient and Hessian of \AtlasInlineMath{f,g} have finite Lipschitz budgets \AtlasInlineMath{L_f,L_g,\rho_f,\rho_g}; \AtlasInlineMath{f(x,\cdot)} is \AtlasInlineMath{C_f}-Lipschitz. Also \AtlasInlineMath{\Phi(x_0)-\inf\Phi\le\Delta} and \AtlasInlineMath{\operatorname{dist}(y_0,Y^*(x_0))\le R_y}. These are population assumptions, not samplewise assumptions.

\needspace{4\baselineskip}\subsection{Oracle interface}
At any \AtlasInlineMath{(x,y)}, choose \AtlasInlineMath{h=f} or \AtlasInlineMath{g} and receive the fresh full-gradient sample \AtlasInlineMath{G_h}. Conditional on history \AtlasInlineMath{\mathcal{H}} and the query, \AtlasInlineMath{\mathbb{E}[G_h\mid\mathcal{H}]=\nabla h} and \AtlasInlineMath{\mathbb{E}[\|G_h-\nabla h\|^2\mid\mathcal{H}]\le\sigma_h^2} hold. Calls use independent fresh samples with no cross-point reuse.

Dimensions, budget, and initial point are public. No direct inner-solution, branch-selection, hypergradient, or second-order interface is supplied.

\needspace{4\baselineskip}\subsection{Output}
Return a single outer point \AtlasInlineMath{\widehat{x}\in\mathbb{R}^{d_x}}, possibly a randomized historical point or an unqueried point. No certificate is required.

\needspace{4\baselineskip}\subsection{Loss or accuracy criterion}
The criterion is \AtlasInlineMath{\|\nabla\Phi(\widehat x)\|}. Differentiability follows from the uniform perturbed-PL model, not from a PL condition on \AtlasInlineMath{g} alone.

\needspace{4\baselineskip}\subsection{Resource accounting}
The budget is \AtlasInlineMath{N=N_{\nabla f}+N_{\nabla g}}. Each sample costs one unit, under a fixed worst-case budget including all initialization and inner-loop calls.

\needspace{4\baselineskip}\subsection{Algorithm views and success guarantees}
Arbitrary randomized adaptive algorithms are allowed, with measurable queries, stopping, and postprocessing. Outputs may be unqueried, randomly selected historical points, or admissible averages. No span or zero-respecting restriction applies. Queries are sequential; each response in a batch is charged. Internal arithmetic, linear algebra, and storage are free, with no separate work, memory, or communication budget.

For every instance, \AtlasInlineMath{\mathbb E\|\nabla\Phi(\widehat x)\|\le\varepsilon}, with expectation over observation and algorithm randomness.

\needspace{7\baselineskip}\section{2. Notation and parameters}
\AtlasInlineMath{x,y}: The upper- and lower-level variables. \textbf{Domain:} \AtlasInlineMath{x\in\mathbb R^{d_x},y\in\mathbb R^{d_y}}.

\AtlasInlineMath{f,g}: The original upper- and lower-level real-valued functions. \textbf{Domain:} \AtlasInlineMath{f,g:\mathbb R^{d_x}\times\mathbb R^{d_y}\to\mathbb R}.

\AtlasInlineMath{\widehat x,x_0,y_0}: Output outer point, public outer initial point, and public lower-level initial point \textbf{Domain:} Each lies in its corresponding variable space.

\AtlasInlineMath{\Phi}: The reduced objective. \textbf{Domain:} Defined in the task specification.

\AtlasInlineMath{Y^*}: Set of lower-level global minimizers \textbf{Domain:} Constraints specified under TASK.

\AtlasInlineMath{\nabla,\|\cdot\|}: Euclidean gradient and vector norm; matrix and higher-derivative norms are the induced Euclidean operator norms

\AtlasInlineMath{\mathbb E,\Pr,\inf,\sup}: Expectation, probability, infimum, and supremum, with aggregation order and randomness specified by the guarantee.

\AtlasInlineMath{\mathcal O,\Omega,\Theta,\widetilde{\mathcal O}}: Asymptotic upper, lower, matching-order, and declared-logarithm upper bounds; fixed parameters are specified in each evidence record

\AtlasInlineMath{c,C}: Positive absolute constants or functions of fixed budgets, with dependencies specified by the corresponding result \textbf{Domain:} \AtlasInlineMath{0<c,C<\infty}.

\AtlasInlineMath{\varepsilon}: Target accuracy. Whether error is squared is determined by the metric and guarantee formulas. \textbf{Domain:} \AtlasInlineMath{\varepsilon>0}.

\AtlasInlineMath{d_x}: Upper-level dimension \textbf{Domain:} Positive integer.

\AtlasInlineMath{d_y}: Lower-level dimension \textbf{Domain:} Positive integer.

\AtlasInlineMath{L_f}: Joint-variable Lipschitz upper bound for the complete population upper-level gradient \textbf{Domain:} \AtlasInlineMath{L_f\ge0}.

\AtlasInlineMath{L_g}: A joint-variable Lipschitz bound on the population lower-level full gradient. \textbf{Domain:} \AtlasInlineMath{L_g\ge\mu>0}.

\AtlasInlineMath{C_f}: A uniform norm bound on the population upper function’s lower-variable gradient block. \textbf{Domain:} \AtlasInlineMath{C_f\ge0}.

\AtlasInlineMath{\|\nabla_yf(x,y)\|\le C_f}

\AtlasInlineMath{\rho_f}: Joint-variable Lipschitz upper bound for the full upper Hessian. \textbf{Domain:} \AtlasInlineMath{\rho_f\ge0}.

\AtlasInlineMath{\rho_g}: Joint-variable Lipschitz upper bound for the full lower Hessian. \textbf{Domain:} \AtlasInlineMath{\rho_g\ge0}.

\AtlasInlineMath{\mu_{\rm PL}}: The common PL constant of all perturbed lower-level objectives. \textbf{Domain:} \AtlasInlineMath{\mu_{\rm PL}>0}.

\AtlasInlineMath{\Delta}: Public upper bound on the initial reduced-objective gap \textbf{Domain:} \AtlasInlineMath{\Delta\ge0}.

\AtlasInlineMath{\Phi(x_0)-\inf_x\Phi(x)\le\Delta}

\AtlasInlineMath{R_y}: A public bound on distance from the initial point to the lower solution set. \textbf{Domain:} \AtlasInlineMath{R_y\ge0}.

\AtlasInlineMath{\operatorname{dist}(y_0,Y^*(x_0))\le R_y}

\AtlasInlineMath{\bar t}: Upper endpoint of the parameter interval in the uniform perturbed PL condition \textbf{Domain:} \AtlasInlineMath{\bar t>0}.

\AtlasInlineMath{\ell}: Common regularity scale used by the uniformly perturbed PL algorithm \textbf{Domain:} \AtlasInlineMath{\ell>0}.

\AtlasInlineMath{\ell=\max\{C_f,L_f,L_g,\rho_g\}}

\AtlasInlineMath{\kappa}: The ratio of the common regularity scale to the uniform PL constant. \textbf{Domain:} \AtlasInlineMath{\kappa>0}.

\AtlasInlineMath{\kappa=\ell/\mu_{\rm PL}}

\AtlasInlineMath{\sigma_f}: Conditional standard-deviation bound for the full stochastic upper-level gradient \textbf{Domain:} \AtlasInlineMath{\sigma_f\ge0}.

\AtlasInlineMath{\sigma_g}: Conditional standard-deviation bound for the full stochastic lower-level gradient \textbf{Domain:} \AtlasInlineMath{\sigma_g\ge0}.

\AtlasInlineMath{G_f,G_g,\mathcal H}: Upper- and lower-level stochastic full-gradient responses and the pre-draw history. \textbf{Domain:} \AtlasInlineMath{G_f,G_g\in\mathbb R^{d_x+d_y}}.

\AtlasInlineMath{G_{f,x},G_{f,y},G_{g,x},G_{g,y}}: Outer and lower blocks of the corresponding full stochastic gradient \textbf{Domain:} Outer-block dimension d\_x and lower-block dimension d\_y.

\AtlasInlineMath{h_t,t}: The perturbed lower-level objective and perturbation parameter. \textbf{Domain:} \AtlasInlineMath{0\le t\le\bar t}.

\AtlasInlineMath{h_t(x,y)=g(x,y)+tf(x,y)}

\AtlasInlineMath{\operatorname{dist}}: Euclidean point-to-set distance

\AtlasInlineMath{\operatorname{dist}(u,S)=\inf_{s\in S}\|u-s\|}

\AtlasInlineMath{N,Q}: This is the total call budget under the stated resource projection. If only one symbol is used, the two names are aliases for that same budget. \textbf{Domain:} RESOURCE specifies which operations are summed; no additional unit cost is introduced.

\needspace{12\baselineskip}\section{3. Complexity bracket}
\needspace{4\baselineskip}\subsection{Compared quantities}
\needspace{4\baselineskip}\subsubsection{Q1 - source scoped metric resource statement}
\textbf{Metric contract.} E||grad Phi(x\_hat)|| <= epsilon; a uniform random outer iterate implements the source average guarantee.

\textbf{Resource contract.} vector: N\_grad\_f; N\_grad\_g; projection: Q=N\_grad\_f+N\_grad\_g; every sample in each batch is charged.; budget: Fixed planned horizon and batch schedule.

This is the native source metric/resource guarantee. Joint precision, block size, runtime and oracle components stay separate; do not infer a scalar optimal query complexity.

\needspace{4\baselineskip}\subsubsection{Q2 - source scoped metric resource statement}
\textbf{Metric contract.} E||grad Phi(x\_hat)|| <= epsilon; a uniform random outer iterate implements the source average guarantee.

\textbf{Resource contract.} vector: N\_grad\_f; N\_grad\_g; projection: Q=N\_grad\_f+N\_grad\_g; every sample in each batch is charged.; budget: Fixed planned horizon and batch schedule.

This is the native source metric/resource guarantee. Joint precision, block size, runtime and oracle components stay separate; do not infer a scalar optimal query complexity.

\needspace{4\baselineskip}\subsubsection{Q3 - source scoped metric resource statement}
\textbf{Metric contract.} E||grad Phi(x\_hat)|| <= epsilon; a uniform random outer iterate implements the source average guarantee.

\textbf{Resource contract.} vector: N\_grad\_f; N\_grad\_g; projection: Q=N\_grad\_f+N\_grad\_g; every sample in each batch is charged.; budget: Fixed planned horizon and batch schedule.

This is the native source metric/resource guarantee. Joint precision, block size, runtime and oracle components stay separate; do not infer a scalar optimal query complexity.

\needspace{4\baselineskip}\subsection{R1 - BM\_BOUND\_pl\_stoch\_det source window}
\textbf{Region.} Assume \AtlasInlineMath{\sigma_f=\sigma_g=0} and take \AtlasInlineMath{\varepsilon} sufficiently small that the source perturbation satisfies \AtlasInlineMath{t_{\rm pen}\le\bar t}. Regularity, variance, initial gap, and distance budgets are fixed. The quantity \AtlasInlineMath{\ell} is the maximum normalized source regularity budget, with \AtlasInlineMath{\kappa=\ell/\mu_{\rm PL}}. The displayed logarithm requires \AtlasInlineMath{\rho_f>0} and \AtlasInlineMath{\rho_f\ell\kappa/\varepsilon\ge e}; a positive public upper bound may be declared even when the actual variation is zero.

\textbf{Quantity.} Q2

\textbf{Resource.} Joint resource budget

\textbf{Status.} Lower bound missing

\textbf{Lower bound.} No admitted bound for this comparison

\textbf{Upper bound.} B1

\needspace{4\baselineskip}\subsection{R2 - BM\_BOUND\_pl\_stoch\_upper\_noise source window}
\textbf{Region.} Assume \AtlasInlineMath{\sigma_f>0,\sigma_g=0} and take \AtlasInlineMath{\varepsilon} sufficiently small that the source perturbation satisfies \AtlasInlineMath{t_{\rm pen}\le\bar t}. Regularity, variance, initial gap, and distance budgets are fixed. The quantity \AtlasInlineMath{\ell} is the maximum normalized source regularity budget, with \AtlasInlineMath{\kappa=\ell/\mu_{\rm PL}}. The displayed logarithm requires \AtlasInlineMath{\rho_f>0} and \AtlasInlineMath{\rho_f\ell\kappa/\varepsilon\ge e}; a positive public upper bound may be declared even when the actual variation is zero.

\textbf{Quantity.} Q3

\textbf{Resource.} Joint resource budget

\textbf{Status.} Lower bound missing

\textbf{Lower bound.} No admitted bound for this comparison

\textbf{Upper bound.} B2

\needspace{4\baselineskip}\subsection{R3 - BM\_BOUND\_pl\_stoch\_both\_noise source window}
\textbf{Region.} Assume \AtlasInlineMath{\sigma_f>0,\sigma_g>0} and take \AtlasInlineMath{\varepsilon} sufficiently small that the source perturbation satisfies \AtlasInlineMath{t_{\rm pen}\le\bar t}. Regularity, variance, initial gap, and distance budgets are fixed. The quantity \AtlasInlineMath{\ell} is the maximum normalized source regularity budget, with \AtlasInlineMath{\kappa=\ell/\mu_{\rm PL}}. The displayed logarithm requires \AtlasInlineMath{\rho_f>0} and \AtlasInlineMath{\rho_f\ell\kappa/\varepsilon\ge e}; a positive public upper bound may be declared even when the actual variation is zero.

\textbf{Quantity.} Q1

\textbf{Resource.} Joint resource budget

\textbf{Status.} Lower bound missing

\textbf{Lower bound.} No admitted bound for this comparison

\textbf{Upper bound.} B3

\needspace{7\baselineskip}\section{4. Bounds and supporting evidence}
The following conditions apply to every bound in this profile, in addition to each bound's stated region and algorithm scope.

\textbf{Oracle mapping.} Selected upper- and lower-level stochastic gradients use fresh samples; same-sample finite differences are unavailable.

\textbf{Resource mapping.} Every gradient response in every batch is charged.

\needspace{4\baselineskip}\subsection{B1 - F2BSA}
\textbf{Admission.} Included in the bracket

\textbf{Direction.} upper

\textbf{Witness or lower-bound scope.} Randomized witness

\AtlasMath{\mathcal{O}\!\left(\ell\kappa^4\varepsilon^{-2}\log\frac{\rho_f\ell\kappa}{\varepsilon}\right)}
\textbf{Sources.} \href{https://proceedings.mlr.press/v247/chen24a/chen24a.pdf}{[1]}

\textbf{Constants.} Hold the stated regularity and initialization budgets fixed. Joint optimality is not claimed for parameters whose dependence is untracked.

\textbf{Source location.} \href{https://proceedings.mlr.press/v247/chen24a/chen24a.pdf}{[1]} pdf pages: 8; 10; 11; 12; theorem or section: Assumption 4.1; Theorems 4.1-4.2; Assumption 4.2; theorem: Assumption 4.1; Theorems 4.1-4.2; Assumption 4.2

\needspace{4\baselineskip}\subsection{B2 - F2BSA}
\textbf{Admission.} Included in the bracket

\textbf{Direction.} upper

\textbf{Witness or lower-bound scope.} Randomized witness

\AtlasMath{\mathcal{O}\!\left(\ell\kappa^5\varepsilon^{-4}\log\frac{\rho_f\ell\kappa}{\varepsilon}\right)}
\textbf{Sources.} \href{https://proceedings.mlr.press/v247/chen24a/chen24a.pdf}{[1]}

\textbf{Constants.} Hold the stated regularity and initialization budgets fixed. Joint optimality is not claimed for parameters whose dependence is untracked.

\textbf{Source location.} \href{https://proceedings.mlr.press/v247/chen24a/chen24a.pdf}{[1]} pdf pages: 8; 10; 11; 12; theorem or section: Assumption 4.1; Theorems 4.1-4.2; Assumption 4.2; theorem: Assumption 4.1; Theorems 4.1-4.2; Assumption 4.2

\needspace{4\baselineskip}\subsection{B3 - F2BSA}
\textbf{Admission.} Included in the bracket

\textbf{Direction.} upper

\textbf{Witness or lower-bound scope.} Randomized witness

\AtlasMath{\mathcal{O}\!\left(\ell^3\kappa^{11}\varepsilon^{-6}\log\frac{\rho_f\ell\kappa}{\varepsilon}\right)}
\textbf{Sources.} \href{https://proceedings.mlr.press/v247/chen24a/chen24a.pdf}{[1]}

\textbf{Constants.} Hold the stated regularity and initialization budgets fixed. Joint optimality is not claimed for parameters whose dependence is untracked.

\textbf{Source location.} \href{https://proceedings.mlr.press/v247/chen24a/chen24a.pdf}{[1]} pdf pages: 8; 10; 11; 12; theorem or section: Assumption 4.1; Theorems 4.1-4.2; Assumption 4.2; theorem: Assumption 4.1; Theorems 4.1-4.2; Assumption 4.2

\needspace{7\baselineskip}\section{References}
[1] Lesi Chen; Jing Xu; Jingzhao Zhang. On Finding Small Hyper-Gradients in Bilevel Optimization: Hardness Results and Improved Analysis. 2024. COLT 2024, PMLR 247, pp.947-980.

\url{https://proceedings.mlr.press/v247/chen24a/chen24a.pdf}

\end{document}
